\chapter{2005年浙江卷（文）}
\section{单选题}
\begin{problem}
函数 \( y = \sin \left(2x + \frac{\pi}{6}\right) \) 的最小正周期是（\quad）

\choices
    {\( \frac{\pi}{2} \)}
    {\(\pi\)}
    {\(2\pi\)}
    {\(4\pi\)}

\begin{answer}
B
\end{answer}
\end{problem}

\begin{problem}
设全集 \(U = \{1, 2, 3, 4, 5, 6, 7\}\)，\(P = \{1, 2, 3, 4, 5\}\)，\(Q = \{3, 4, 5, 6, 7\}\)，则 \(P \cap (\complement_U Q) =\)（\quad）

\choices
    {\(\{1,2\}\)}
    {\(\{3,4,5\}\)}
    {\(\{1,2,6,7\}\)}
    {\(\{1,2,3,4,5\}\)}

\begin{answer}
A
\end{answer}
\end{problem}

\begin{problem}
点 \((1, -1)\) 到直线 \(x - y + 1 = 0\) 的距离是（\quad）

\choices
    {\( \frac{1}{2} \)}
    {\( \frac{3}{2} \)}
    {\( \frac{\sqrt{2}}{2} \)}
    {\(\frac{3\sqrt{2}}{2}\)}

\begin{answer}
D
\end{answer}
\end{problem}

\begin{problem}
设 \( f(x) = |x - 1| - |x| \)，则 \( f\left[f\left(\frac{1}{2}\right)\right] = \)（\quad）

\choices
    {\(-\frac{1}{2}\)}
    {\(0\)}
    {\( \frac{1}{2} \)}
    {\(1\)}

\begin{answer}
D
\end{answer}
\end{problem}

\begin{problem}
在 \((1 - x)^{5} - (1 - x)^{6}\) 的展开式中，含 \(x^{3}\) 的项的系数是（\quad）

\choices
    {\(-5\)}
    {\(5\)}
    {\(-10\)}
    {\(10\)}

\begin{answer}
D
\end{answer}
\end{problem}

\begin{problem}
从存放号码分别为 \(1\)，\(2\)，\(\cdots\)，\(10\) 的卡片的盒子中，有放回地取 \(100\) 次，每次取一张卡片并记下号码，统计结果如下：

\begin{center}
\begin{tabular}{|c|c c c c c c c c c c|}
\hline
卡片号码 & 1 & 2 & 3 & 4 & 5 & 6 & 7 & 8 & 9 & 10 \\
\hline
取到的次数 & 13 & 8 & 5 & 7 & 6 & 13 & 18 & 10 & 11 & 9 \\
\hline
\end{tabular}
\end{center}

则取到号码为奇数的频率是（\quad）

\choices
    {\(0.53\)}
    {\(0.5\)}
    {\(0.47\)}
    {\(0.37\)}

\begin{answer}
A
\end{answer}
\end{problem}

\begin{problem}
设 \(\alpha\)，\(\beta\) 为两个不同的平面，\(l\)，\(m\) 为两条不同的直线，且 \(l \subset \alpha\)，\(m \subset \beta\)，有如下的两个命题：
\begin{circlelist}
\item 若 \(\alpha \parallel \beta\)，则 \(l \parallel m\)；
\item 若 \(l \perp m\)，则 \(\alpha \perp \beta\).
\end{circlelist}
那么（\quad）

\choices
    {\circled{1}是真命题，\circled{2}是假命题}
    {\circled{1}是假命题，\circled{2}是真命题}
    {\circled{1}\circled{2}都是真命题}
    {\circled{1}\circled{2}都是假命题}

\begin{answer}
D
\end{answer}
\end{problem}

\begin{problem}
已知向量 \(\bs{a} = (x - 5, 3)\)，\(\bs{b} = (2, x)\)，且 \(\bs{a} \perp \bs{b}\)，则由 \(x\) 值的构成的集合是（\quad）

\choices
    {\( \{2,3\} \)}
    {\(\{-1,6\}\)}
    {\( \{2\} \)}
    {\( \{6\} \)}

\begin{answer}
C
\end{answer}
\end{problem}

\begin{problem}
函数 \( y = ax^{2} + 1 \) 的图象与直线 \( y = x \) 相切，则 \( a = \)（\quad）

\choices
    {\( \frac{1}{8} \)}
    {\( \frac{1}{4} \)}
    {\( \frac{1}{2} \)}
    {\(1\)}

\begin{answer}
B
\end{answer}
\end{problem}

\begin{problem}
设集合 \(A = \{(x,y)\mid x\)，\(y\)，\(1-x-y\) 是三角形的三边长\(\}\)，则 \(A\) 所表示的平面区域（不含边界的阴影部分）是（\quad）

\choices
    {\choicebitmap[width=0.15\paperwidth]{img/2005/zhejiang_liberal/q10_optA.png}}
    {\choicebitmap[width=0.15\paperwidth]{img/2005/zhejiang_liberal/q10_optB.png}}
    {\choicebitmap[width=0.15\paperwidth]{img/2005/zhejiang_liberal/q10_optC.png}}
    {\choicebitmap[width=0.15\paperwidth]{img/2005/zhejiang_liberal/q10_optD.png}}

\begin{answer}
A
\end{answer}
\end{problem}

\section{填空题}
\begin{problem}
函数 \( y = \frac{x}{x + 2} \)（\(x \in \R\)，且 \(x \neq -2\)）的反函数是\fillinblank{}.

\begin{answer}
\(y=\frac{2x}{1-x}\)（\(x\in\R\)，且 \(x\neq1\)）
\end{answer}
\end{problem}

\begin{problem}
设 \(M\)，\(N\) 是直角梯形 \(ABCD\) 两腰的中点，\(DE \perp AB\) 于 \(E\) （如图）. 现将 \(\triangle ADE\) 沿 \(DE\) 折起，使二面角 \(A - DE - B\) 为 \(45^\circ\)，此时点 \(A\) 在平面 \(BCDE\) 内的射影恰为点 \(B\)，则 \(M\)，\(N\) 的连线与 \(AE\) 所成角的大小等于\fillinblank{}.

\examfiguregroup[float=inline]{img/2005/zhejiang_liberal/q12_fig1.png}{%
    \bitmapinclude[max width=0.55\linewidth]{img/2005/zhejiang_liberal/q12_fig1.png}\hfill
    \bitmapinclude[max width=0.35\linewidth]{img/2005/zhejiang_liberal/q12_fig2.png}%
}

\begin{answer}
\(90^\circ\)
\end{answer}
\end{problem}

\begin{problem}
过双曲线 \( \frac{x^{2}}{a^{2}} - \frac{y^{2}}{b^{2}} = 1 \)（\(a > 0\)，\(b > 0\)）的左焦点且垂直于 \(x\) 轴的直线与双曲线相交于 \(M\)，\(N\) 两点，以 \(MN\) 为直径的圆恰好过双曲线的右顶点，则双曲线的离心率等于\fillinblank{}.

\begin{answer}
\(2\)
\end{answer}
\end{problem}

\begin{problem}
从集合 \(\{P, Q, R, S\}\) 与 \(\{0, 1, 2, 3, 4, 5, 6, 7, 8, 9\}\) 中各任取 2 个元素排成一排（字母和数字均不能重复）. 每排中字母 \(Q\) 和数字 \(0\) 至多只能出现一个的不同排法种数是\fillinblank{}. （用数字作答）

\begin{answer}
\(5832\)
\end{answer}
\end{problem}

\section{解答题}
\begin{problem}
已知函数 \( f(x) = 2\sin x\cos x + \cos 2x \).
\begin{enumerate}
\item 求 \( f\left(\frac{\pi}{4}\right) \) 的值；
\item 设 \(\alpha \in (0, \pi)\)，\(f\left(\frac{\alpha}{2}\right) = \frac{\sqrt{2}}{2}\)，求 \(\sin \alpha\) 的值.
\end{enumerate}

\begin{solution}
\begin{enumerate}
\item 由二倍角公式，\(f(x)=2\sin x\cos x+\cos 2x=\sin 2x+\cos 2x\)，所以 \(f\left(\frac{\pi}{4}\right)=\sin\frac{\pi}{2}+\cos\frac{\pi}{2}=1\).

\item 由二倍角公式，\(f\left(\frac{\alpha}{2}\right)=2\sin\frac{\alpha}{2}\cos\frac{\alpha}{2}+\cos\alpha=\sin\alpha+\cos\alpha=\frac{\sqrt{2}}{2}\). 于是 \(\cos\alpha=\frac{\sqrt{2}}{2}-\sin\alpha\). 代入 \(\sin^2\alpha+\cos^2\alpha=1\)，得 \(\sin^2\alpha+\left(\frac{\sqrt{2}}{2}-\sin\alpha\right)^2=1\)，即 \(2\sin^2\alpha-\sqrt{2}\sin\alpha-\frac{1}{2}=0\). 解得 \(\sin\alpha=\frac{\sqrt{2}+\sqrt{6}}{4}\) 或 \(\sin\alpha=\frac{\sqrt{2}-\sqrt{6}}{4}\). 因为 \(\alpha\in(0,\pi)\)，所以 \(\sin\alpha>0\)，故 \(\sin\alpha=\frac{\sqrt{2}+\sqrt{6}}{4}\).
\end{enumerate}
\end{solution}
\end{problem}

\begin{problem}
已知实数 \(a\)，\(b\)，\(c\) 成等差数列，\(a + 1\)，\(b + 1\)，\(c + 4\) 成等比数列. 且 \(a + b + c = 15\)，求 \(a\)，\(b\)，\(c\).

\begin{solution}
因为 \(a\)，\(b\)，\(c\) 成等差数列，所以 \(a+c=2b\). 又 \(a+b+c=15\)，从而 \(3b=15\)，即 \(b=5\).

设等差数列的公差为 \(d\)，则 \(a=5-d\)，\(c=5+d\). 由 \(a+1\)，\(b+1\)，\(c+4\) 成等比数列，得 \((b+1)^2=(a+1)(c+4)\)，即 \(36=(6-d)(9+d)=54-3d-d^2\). 所以 \(d^2+3d-18=(d-3)(d+6)=0\)，故 \(d=3\) 或 \(d=-6\).

当 \(d=3\) 时，\((a,b,c)=(2,5,8)\)；当 \(d=-6\) 时，\((a,b,c)=(11,5,-1)\). 因此 \((a,b,c)=(2,5,8)\) 或 \((a,b,c)=(11,5,-1)\).
\end{solution}
\end{problem}

\begin{problem}
袋子 \(A\) 和 \(B\) 中装有若干个均匀的红球和白球，从 \(A\) 中摸出一个红球的概率是 \(\frac{1}{3}\)，从 \(B\) 中摸出一个红球的概率为 \(p\).
\begin{enumerate}
\item 从 \(A\) 中有放回地摸球，每次摸出一个，共摸 5 次.
\begin{enumerate}
\item 恰好有 3 次摸到红球的概率；
\item 第一次，第三次，第五次摸到红球的概率；
\end{enumerate}
\item 若 \(A\)，\(B\) 两个袋子中的球数之比为 \(1:2\)，将 \(A\)，\(B\) 中的球装在一起后，从中摸出一个红球的概率是 \( \frac{2}{5} \)，求 \(p\) 的值.
\end{enumerate}

\begin{solution}
\begin{enumerate}
\item 从袋子 \(A\) 中有放回地摸球，所以每次摸到红球的概率均为 \(\frac{1}{3}\)，各次摸球相互独立.
\begin{enumerate}
\item 恰好有 \(3\) 次摸到红球的情况数为 \(\mathrm C_5^3\)，所以所求概率为
\[
\mathrm C_5^3\left(\frac{1}{3}\right)^3\left(1-\frac{1}{3}\right)^2
=10\cdot\frac{1}{27}\cdot\frac{4}{9}
=\frac{40}{243}.
\]

\item 第一次、第三次、第五次摸到红球，而第二次、第四次摸到何种球均不受限制，所以所求概率为
\(\frac{1}{3}\cdot\frac{1}{3}\cdot\frac{1}{3}=\frac{1}{27}\).
\end{enumerate}

\item 设袋子 \(A\) 中有 \(n\) 个球，则袋子 \(B\) 中有 \(2n\) 个球. 袋子 \(A\) 中红球数为 \(\frac{n}{3}\)，袋子 \(B\) 中红球数为 \(2np\). 装在一起后，红球概率为 \(\frac{\frac{n}{3}+2np}{3n}=\frac{2}{5}\). 约去 \(n\)，得 \(\frac{1}{9}+\frac{2p}{3}=\frac{2}{5}\)，所以 \(\frac{2p}{3}=\frac{13}{45}\)，从而 \(p=\frac{13}{30}\).
\end{enumerate}
\end{solution}
\end{problem}

\begin{problem}
如图，在三棱锥 \(P - ABC\) 中，\(AB \perp BC\)，\(AB = BC = \frac{1}{2} PA\)，点 \(O\)，\(D\) 分别是 \(AC\)，\(PC\) 的中点，\(OP \perp\) 底面 \(ABC\).
\begin{enumerate}
\item 求证： \( OD \parallel \) 平面 \(PAB\)；
\item 求直线 \(PA\) 与平面 \(PBC\) 所成角的大小.
\end{enumerate}

\bitmapfigure[max width=0.36\linewidth]{img/2005/zhejiang_liberal/q18_fig1.png}

\begin{solution}
\begin{enumerate}
\item 在三角形 \(APC\) 中，\(O\)，\(D\) 分别是 \(AC\)，\(PC\) 的中点，所以由中位线定理，
\(OD\parallel AP\). 又因为 \(AP\subset\) 平面 \(PAB\). 平面 \(PAB\) 与底面 \(ABC\) 的交线为 \(AB\)，而点 \(O\) 是斜边 \(AC\) 的中点，不在 \(AB\) 上，所以 \(O\notin\) 平面 \(PAB\). 因此 \(OD\parallel\) 平面 \(PAB\).

\item 设 \(AB=BC=s\)，其中 \(s>0\). 建立空间直角坐标系，使底面 \(ABC\) 在 \(z=0\) 平面内，取 \(B(0,0,0)\)，\(A(s,0,0)\)，\(C(0,s,0)\). 则 \(O\) 为 \(AC\) 的中点，且 \(OP\perp\) 底面，所以可设 \(O\left(\frac{s}{2},\frac{s}{2},0\right)\)，\(P\left(\frac{s}{2},\frac{s}{2},h\right)\)，其中 \(h>0\). 由 \(PA=2s\)，有 \(\left(\frac{s}{2}\right)^2+\left(\frac{s}{2}\right)^2+h^2=4s^2\)，所以 \(h^2=\frac{7s^2}{2}\).

取直线 \(AP\) 的方向向量 \(\bs{v}=\overrightarrow{AP}=\left(-\frac{s}{2},\frac{s}{2},h\right)\)，平面 \(PBC\) 内的两个方向向量为 \(\overrightarrow{BP}=\left(\frac{s}{2},\frac{s}{2},h\right)\) 和 \(\overrightarrow{BC}=(0,s,0)\). 因而平面 \(PBC\) 的一个法向量可取 \(\bs{n}=\overrightarrow{BP}\times\overrightarrow{BC}=\left(-hs,0,\frac{s^2}{2}\right)\). 设直线 \(PA\) 与平面 \(PBC\) 所成的角为 \(\theta\)，则 \(\sin\theta=\frac{\left|\bs{v}\cdot\bs{n}\right|}{\left|\bs{v}\right|\left|\bs{n}\right|}\). 其中 \(\left|\bs{v}\right|=PA=2s\)，\(\bs{v}\cdot\bs{n}=hs^2\)，并且 \(\left|\bs{n}\right|=s\sqrt{h^2+\frac{s^2}{4}}=\frac{s^2\sqrt{15}}{2}\). 所以
\[
\sin\theta=\frac{hs^2}{2s\cdot\frac{s^2\sqrt{15}}{2}}=\frac{h}{s\sqrt{15}}=\sqrt{\frac{7}{30}}.
\]
因此 \(\theta=\arcsin\sqrt{\frac{7}{30}}\).
\end{enumerate}
\end{solution}
\end{problem}

\begin{problem}
如图，已知椭圆的中心在坐标原点，焦点 \(F_{1}\)，\(F_{2}\) 在 \(x\) 轴上，长轴 \(A_{1}A_{2}\) 的长为 4，左准线 \(l\) 与 \(x\) 轴的交点为 \(M\)，\(|MA_{1}|: |A_{1}F_{1}| = 2: 1\).
\begin{enumerate}
\item 求椭圆的方程；
\item 若点 \(P\) 为 \(l\) 上的动点，求 \(\angle F_1PF_2\) 的最大值.
\end{enumerate}

\bitmapfigure[max width=0.44\linewidth]{img/2005/zhejiang_liberal/q19_fig1.png}

\begin{solution}
\begin{enumerate}
\item 设椭圆的半长轴为 \(a\)，半焦距为 \(c\)，半短轴为 \(b\)，则 \(2a=4\)，\(a^2=b^2+c^2\)，所以 \(a=2\). 左准线方程为 \(x=-\frac{a^2}{c}=-\frac{4}{c}\)，故 \(M\left(-\frac{4}{c},0\right)\)，\(A_{1}(-2,0)\)，\(F_{1}(-c,0)\). 由题意，\(0<c<2\)，且 \(\frac{|MA_{1}|}{|A_{1}F_{1}|}=\frac{\frac{4}{c}-2}{2-c}=2\). 化简得 \(\frac{2(2-c)}{c(2-c)}=2\)，所以 \(c=1\). 于是 \(b^2=a^2-c^2=4-1=3\)，椭圆的方程为 \(\frac{x^2}{4}+\frac{y^2}{3}=1\).

\item 由第（1）问，\(F_{1}(-1,0)\)，\(F_{2}(1,0)\)，且直线 \(l\) 的方程为 \(x=-4\). 设动点 \(P\left(-4,t\right)\)，令 \(u=|t|\). 此时 \(\overrightarrow{PF_{1}}=(3,-t)\)，\(\overrightarrow{PF_{2}}=(5,-t)\)，两向量的数量积为 \(t^2+15>0\)，故 \(\theta=\angle F_{1}PF_{2}\) 为锐角，并且
\[
\tan\theta=\frac{\left|3(-t)-(-t)5\right|}{3\cdot5+(-t)(-t)}=\frac{2u}{u^2+15}.
\]
由基本不等式 \(u^2+15\geqslant2\sqrt{15}\,u\)，得 \(\tan\theta=\frac{2u}{u^2+15}\leqslant\frac{1}{\sqrt{15}}\). 当且仅当 \(u=\sqrt{15}\)，即 \(P\left(-4,\sqrt{15}\right)\) 或 \(P\left(-4,-\sqrt{15}\right)\) 时等号成立. 因为 \(\theta\) 为锐角，所以其最大值为 \(\arctan\frac{1}{\sqrt{15}}\).
\end{enumerate}
\end{solution}
\end{problem}

\begin{problem}
已知函数 \( f(x) \) 和 \( g(x) \) 的图象关于原点对称，且 \( f(x) = x^2 + 2x \).
\begin{enumerate}
\item 求函数 \( g(x) \) 的解析式；
\item 解不等式 \( g(x) \geqslant f(x) - |x - 1| \)；
\item 若 \(h(x) = g(x) - \lambda f(x) + 1\) 在 \([-1, 1]\) 上是增函数，求实数 \(\lambda\) 的取值范围.
\end{enumerate}

\begin{solution}
\begin{enumerate}
\item 图象关于原点对称，所以点 \((x,g(x))\) 关于原点的对称点 \((-x,-g(x))\) 在函数 \(f\) 的图象上，即 \(f(-x)=-g(x)\). 因此 \(g(x)=-f(-x)=-\left((-x)^2+2(-x)\right)=-x^2+2x\).

\item 由 \(g(x)\geqslant f(x)-|x-1|\)，得 \(-x^2+2x\geqslant x^2+2x-|x-1|\)，即 \(|x-1|\geqslant2x^2\). 当 \(x<1\) 时，\(|x-1|=1-x\)，所以 \(2x^2+x-1\leqslant0\)，即 \((2x-1)(x+1)\leqslant0\)，解得 \(-1\leqslant x\leqslant\frac{1}{2}\).

当 \(x\geqslant1\) 时，\(|x-1|=x-1\)，所以 \(2x^2-x+1\leqslant0\). 该二次式的判别式为 \(1-8=-7<0\)，且二次项系数为正，故无实数解. 因此原不等式的解集为 \(\left[-1,\frac{1}{2}\right]\).

\item \(h(x)=g(x)-\lambda f(x)+1=-(1+\lambda)x^2+2(1-\lambda)x+1\)，所以 \(h'(x)=-2(1+\lambda)x+2(1-\lambda)\). 这是关于 \(x\) 的一次函数，且 \(h'(-1)=4\)，\(h'(1)=-4\lambda\). 要使 \(h(x)\) 在 \([-1,1]\) 上为增函数，需且只需 \(h'(x)\geqslant0\) 对区间内所有 \(x\) 成立. 由于 \(h'\) 为一次函数，这等价于两个端点处均非负，故 \(-4\lambda\geqslant0\). 因此实数 \(\lambda\) 的取值范围为 \(\lambda\leqslant0\).
\end{enumerate}
\end{solution}
\end{problem}
